% preface_sections5_9_draft.tex — Sections 5-9 of the Restored Preface
% To be integrated into preface.tex
% Macros assumed from main.tex: \cA, \barB, \gAmod, \MC, \cZ
\providecommand{\cA}{\mathcal{A}}
\providecommand{\barB}{\overline{B}}
\providecommand{\gAmod}{\mathfrak{g}_{\cA}^{\mathrm{mod}}}
\providecommand{\MC}{\text{MC}}
\providecommand{\cM}{\mathcal{M}}
\providecommand{\cC}{\mathcal{C}}
\providecommand{\cO}{\mathcal{O}}
\providecommand{\cP}{\mathcal{P}}
\providecommand{\Eone}{E_1}
\providecommand{\Convinf}{\mathrm{Conv}_\infty}
\providecommand{\Convstr}{\mathrm{Conv}_{\mathrm{str}}}
\providecommand{\fg}{\mathfrak{g}}
\providecommand{\Defcyc}{\mathrm{Def}_{\mathrm{cyc}}}

% ====================================================================
\section*{5.\quad The modular tangent complex and Chern--Weil theory}
% ====================================================================

\subsection*{5.1.\enspace Three characteristic projections}

Three successive projections of $\Theta_\cA$ capture progressively
finer structure.

\medskip

\noindent\textbf{The scalar curvature $\kappa(\cA)$.}\enspace
The degree-two, genus-zero trace:
\[
\kappa(\cA)
\;=\;
\operatorname{tr}(\Theta_\cA)_{2,0}
\;=\;
\mathrm{av}\bigl(r(z)\bigr)
\;\in\;\Bbbk.
\]
This is the first Chern class of the modular bar bundle:
$\operatorname{obs}_g(\cA)=\kappa(\cA)\cdot\lambda_g$ for
uniform-weight algebras at all genera, and unconditionally
at genus~$1$. On
the $E_1$ side, $\kappa$ is the $\Sigma_2$-coinvariant of
the $R$-matrix $r(z)$; averaging collapses the $z$-dependent
profile to a number. Explicit values:
\begin{center}
\renewcommand{\arraystretch}{1.2}
\begin{tabular}{lccl}
\textbf{Algebra $\cA$}
 & $\kappa(\cA)$
 & $\kappa(\cA^!)$
 & \textbf{Sum rule}\\
\hline
Heisenberg $\cH_k$ & $k$ & $-k$ & $0$\\[2pt]
Affine $\widehat{\fg}_k$
 & $\dfrac{(k+h^\vee)\dim\fg}{2h^\vee}$
 & $-\dfrac{(k+h^\vee)\dim\fg}{2h^\vee}$ & $0$\\[6pt]
Virasoro $\mathrm{Vir}_c$ & $c/2$ & $(26-c)/2$ & $13$\\[2pt]
$\beta\gamma$ system & $1$ & $-1$ & $0$\\[2pt]
Lattice $V_\Lambda$ & $\operatorname{rank}(\Lambda)$
 & $-\operatorname{rank}(\Lambda)$ & $0$\\[2pt]
$\mathcal W_N(\widehat{\mathfrak{sl}}_N)$
 & $c\cdot(H_N{-}1)$
 & $c'\cdot(H_N{-}1)$
 & $K_N\cdot(H_N{-}1)$\\[6pt]
Free fermion $\psi$ & $1/4$ & $-1/4$ & $0$
\end{tabular}
\end{center}
The complementarity sum $\kappa(\cA)+\kappa(\cA^!)=0$ holds for
Kac--Moody and free-field algebras. For $\mathcal W$-algebras the
sum is a nonzero constant: $13$ for Virasoro, $250/3$
for~$\mathcal W_3$.

\medskip

\noindent\textbf{The spectral discriminant $\Delta_\cA(x)$.}\enspace
The Koszul dual Hilbert series
$h_{\cA^!}(q)=\sum_{n\ge 0}\dim H^n(\barB(\cA))\,q^n$ satisfies a
holonomic recurrence of finite order~$d$. The transfer matrix
$T_{\mathrm{br},\cA}$ of this recurrence acts on a
finite-rank perfect complex $V^{\mathrm{br}}_\cA$, and
\[
\Delta_\cA(x)
\;:=\;
\operatorname{sdet}\!\bigl(1-x\,T_{\mathrm{br},\cA}\bigr).
\]
The eigenvalues of $T_{\mathrm{br}}$ are the reciprocal branch
points of the analytic continuation of~$h_{\cA^!}$. Three
structural properties: (i)~self-duality
$\Delta_{\cA^!}=\Delta_\cA$; (ii)~functoriality under exact
factorization; (iii)~spectral radius
$\rho(T_{\mathrm{br}})=\dim\fg$ for affine algebras at
generic level.

For $\widehat{\mathfrak{sl}}_3$: a degree-$3$ recurrence with
eigenvalues $8$, $(3\pm\sqrt{13})/2$, giving
$\Delta_{\widehat{\mathfrak{sl}}_3}(x)
=(1-8x)(1-3x-x^2)$. The dominant eigenvalue
$8=\dim(\mathfrak{sl}_3)$ controls exponential growth. For
Heisenberg: $\Delta_{\cH}(x)=1$.

In the Chern--Weil dictionary, the logarithmic derivative
$\operatorname{ch}_r(\cA)
=\frac{1}{r}\operatorname{tr}(T_{\mathrm{br}}^r)$
recovers the power-sum symmetric functions of the eigenvalues:
$\Delta_\cA$ plays the role of the Chern character.

\medskip

\noindent\textbf{The quartic shadow
$\mathfrak R^{\mathrm{mod}}_{4}(\cA)$.}\enspace
The first \emph{nonlinear} characteristic class: the first
invariant not recoverable from~$\kappa$ and~$\Delta_\cA$. Its
genus-zero, four-point specialisation is the quartic contact
invariant $Q^{\mathrm{contact}}(\cA)$:
\[
Q^{\mathrm{contact}}_{\mathrm{Vir}}
\;=\;
\frac{10}{c(5c+22)},
\qquad
Q^{\mathrm{contact}}_{\cH}=0,
\qquad
Q^{\mathrm{contact}}_{\beta\gamma}
\;\neq\; 0.
\]

\subsection*{5.2.\enspace The Chern--Weil dictionary}

The modular bar bundle
$\mathrm{B}^{\log}_{\mathrm{mod}}(\cA)\to\overline{\cM}_{g,n}$
is a principal bundle for the gauge group of the convolution
algebra; $\Theta_\cA$ is its flat connection; the shadow
tower is its Chern--Weil theory.

\begin{center}
\renewcommand{\arraystretch}{1.25}
\begin{tabular}{ll}
\textbf{Classical Chern--Weil}
 & \textbf{Modular factorization}\\
\hline
Principal $G$-bundle $P\to M$
 & Modular bar bundle
 $\mathrm{B}^{\log}_{\mathrm{mod}}(\cA)\to\overline{\cM}_{g,n}$\\[2pt]
Connection $1$-form $\omega\in\Omega^1(P,\fg)$
 & MC element $\Theta_\cA\in\mathrm{MC}(\gAmod)$\\[2pt]
Curvature $\Omega=d\omega+\frac{1}{2}[\omega,\omega]$
 & MC equation $D\Theta_\cA+\tfrac12[\Theta_\cA,\Theta_\cA]=0$\\[2pt]
Flatness $\Omega=0$
 & Automatic (bar-intrinsic)\\[2pt]
First Chern class $c_1$
 & Scalar curvature $\kappa(\cA)=\mathrm{av}(r(z))$\\[2pt]
Chern character
 & Spectral discriminant
 $\Delta_\cA(t)=\operatorname{sdet}(1-tT_{\mathrm{br}})$\\[2pt]
Higher Chern classes
 & Shadows $\mathrm{Sh}_r$ for $r\ge 4$\\[2pt]
Weil homomorphism
 & Shadow algebra
 $\cA^{\mathrm{sh}}=H_\bullet(\Defcyc^{\mathrm{mod}})$\\[2pt]
Lagrangian splitting $T^*G=\fg\oplus\fg^*$
 & Complementarity $Q_g(\cA)\oplus Q_g(\cA^!)$\\[2pt]
Holonomy
 & Factorisation homology $\int_\Sigma\cA$\\[2pt]
Gauge equivalence
 & MC gauge in
 $\mathrm{MC}_\bullet(\gAmod)$\\[2pt]
Transgression
 & Shadow obstruction tower (\S6)
\end{tabular}
\end{center}

\subsection*{5.3.\enspace The shadow CohFT}

The MC equation
$D\Theta_\cA+\tfrac{1}{2}[\Theta_\cA,\Theta_\cA]=0$ holds in
the modular convolution algebra. Projecting to genus~$g$ and
degree~$n$ produces a class
$\tau_{g,n}(\cA)\in R^*(\overline{\cM}_{g,n+1})$ in the
tautological ring, and the MC equation at $(g,n)$ descends to
a tautological relation.

The mechanism: the bracket $[\Theta_\cA,\Theta_\cA]$
decomposes according to the boundary of
$\overline{\cM}_{g,n}$. Separating degenerations (a curve
splitting into two components joined by a node) produce
$[\Theta,\Theta]|_{\mathrm{sep}}$; nonseparating degenerations
(a curve acquiring a self-node, genus rising by~$1$) produce
$\hbar\,\Delta(\Theta)|_{\mathrm{nsep}}$. Each boundary
stratum contributes one term; the MC equation forces their sum
to be $D$-exact. This is a tautological relation.

Two instances fix the pattern. At genus~$0$, degree~$3$: the
separating term produces the WDVV equation (three channels of
the four-point function on $\overline{\cM}_{0,4}\cong\mathbb
P^1$, one for each boundary point). At degree~$2$,
genus~$g\ge 2$: the nonseparating loop contraction produces
the Mumford $\lambda$-class clutching relation, with the
coefficient~$\kappa(\cA)$ from the bar curvature. The
complementarity propagator $P_\cA$ is the Givental $R$-matrix;
when a flat unit exists, Teleman reconstruction recovers the
full collection $\{\tau_{g,n}\}$ at all $(g,n)$ from genus-$0$
data. The recursion that computes $\tau_{g,n}$ from
$\tau_{g-1,n+2}$ and $\tau_{g_1,n_1}\otimes\tau_{g_2,n_2}$ is
Eynard--Orantin topological recursion: the MC equation applied
to the shadow obstruction tower.

The classes $\tau_{g,n}(\cA)$ assemble into a cohomological
field theory (Theorem~\ref{thm:shadow-cohft}): $D^2=0$ on the
log-FM compactification implies equivariance and splitting.
The flat identity holds conditionally, when the vacuum
$|0\rangle$ lies in the generating space~$V$. For class-M
algebras, the infinite shadow tower generates tautological
classes at every genus; at genus~$2$--$3$, the relations for
$\mathrm{Vir}_c$ at generic~$c$ produce elements of the Pixton
ideal in $R^*(\overline{\cM}_{g,n})$: Pixton's relations arise
as shadows of the Virasoro MC equation
(Theorem~\ref{thm:pixton-from-mc-semisimple}).

\subsection*{5.4.\enspace The genus spectral sequence}

The genus filtration $F^g\gAmod:=\bigoplus_{g'\ge g}
(\gAmod)_{g',*,*}$ yields a convergent spectral sequence
distinct from the PBW spectral sequence. PBW computes
\emph{what} bar cohomology is (within a fixed genus);
the genus spectral sequence determines \emph{whether} the MC
element extends across genera.

The $E_1$ page isolates tree-level ($p=0$), one-loop ($p=1$),
and genus-two shell ($p=2$) data. The differential
$d_1\colon E_1^{0,*}\to E_1^{1,*}$ is the genus-one
obstruction $\mathrm{Ob}_1$: the obstruction to lifting a
tree-level MC element to one-loop. Convergence follows from
completeness of the depth filtration and pronilpotence of the
weight filtration: the $\varprojlim^1$ term vanishes by
Mittag-Leffler.

For Heisenberg, the spectral sequence degenerates at~$E_1$.
For Virasoro, $d_1\neq 0$ and nontrivial differentials persist
at every page.

\subsection*{5.5.\enspace Homotopy invariance}

The shadow algebra
$\cA^{\mathrm{sh}}=H_\bullet(\Defcyc^{\mathrm{mod}}(\cA))$ is a
homotopy invariant. For any $\infty_\alpha$-quasi-isomorphism
$f\colon\cA\xrightarrow{\sim}\cA'$, the induced map on MC
$\infty$-groupoids is a weak homotopy equivalence. The scalar
$\kappa(\cA)$, the spectral discriminant $\Delta_\cA(t)$, and
the full shadow obstruction tower are quasi-isomorphism
invariants. The MC element $\Theta_\cA$ is invariant up to
gauge: $f_*(\Theta_\cA)$ and $\Theta_{\cA'}$ are
gauge-equivalent.

The assignment
$\cA\mapsto(\gAmod,\Theta_\cA)$ is a functor from chiral
algebras (with $\infty_\alpha$-quasi-isomorphisms inverted) to MC
moduli problems: modular Koszul duality is a derived invariant.


% ====================================================================
\section*{6.\quad The shadow obstruction tower}
% ====================================================================

\subsection*{6.1.\enspace The extension tower}

Define the total weight $w(g,r,d):=2g-2+r+d$. The total weight
filtration $F^N\mathfrak{g}^{\mathrm{amb}}_\cA$ is exhaustive,
separated, and complete. The \emph{modular extension tower}:
\[
\cE_\cA(N)
:=\operatorname{MC}\bigl(\mathfrak g^{\mathrm{amb}}_\cA/F^{N+1}\bigr),
\qquad N\ge 1.
\]
The bar-intrinsic construction provides a compatible point
$\Theta_\cA\in\varprojlim_{N}\cE_\cA(N)$ unconditionally.

\subsection*{6.2.\enspace Obstruction classes}

The truncated MC elements
$\Theta_\cA^{\le 2},\;\Theta_\cA^{\le 3},
\;\Theta_\cA^{\le 4},\;\ldots$
satisfy the MC equation up to a remainder:
\[
D\Theta^{\le r}_\cA
+\tfrac12[\Theta^{\le r}_\cA,\Theta^{\le r}_\cA]
\;=\;
-o_{r+1}(\cA),
\qquad
o_{r+1}(\cA)
\;\in\;
Z^2\!\left(\frac{F^{r+1}\mathfrak g^{\mathrm{amb}}_\cA}
{F^{r+2}\mathfrak g^{\mathrm{amb}}_\cA},\;d_{\le r}\right).
\]
The bar-intrinsic construction guarantees that
$[o_{r+1}]=0$; the shadow tower records the gauge choices.

\subsection*{6.3.\enspace Cubic gauge triviality and the canonical quartic}

When $H^1(F^3\mathfrak g/F^4\mathfrak g,\,d_2)=0$, the cubic
MC term $\Theta^{[3]}_\cA$ is gauge-trivial. After the gauge
transformation, the quartic class
$[\mathfrak Q(\cA)]\in H^2(F^4\mathfrak g/F^5\mathfrak g,d_2)$
is \emph{canonical}: independent of retraction data, cubic gauge,
and orbit representative. This is the first characteristic class
not polynomial in~$\kappa$. The quartic contact invariant:
\[
Q^{\mathrm{contact}}_{\mathrm{Vir}}
\;=\;
\frac{10}{c(5c+22)}.
\]
The denominator $c(5c+22)$ has physical meaning: $c=0$ is the
logarithmic threshold; $c=-22/5$ is the minimal model
accumulation point $\mathcal M(2,5)$.

The clutching law for the quartic class follows from Mok's log
FM degeneration formula: on a degeneration $C\leadsto C_0$, the
quartic restricts with a planted-forest correction that is
genuinely logarithmic, with no classical counterpart.

\subsection*{6.4.\enspace The algebraicity theorem}

On each primary deformation line~$L$, the shadow generating
function $H(t)=\sum_{r\ge 2}rS_r\,t^r$ satisfies
\[
H(t)^2\;=\;t^4\,Q_L(t),
\qquad
Q_L(t)=4\kappa^2+12\kappa\alpha\,t+(9\alpha^2+16\kappa S_4)t^2,
\]
a quadratic polynomial in three genus-$0$ invariants
$(\kappa,\alpha,S_4)$
(Theorem~\ref{thm:riccati-algebraicity}). The infinite tower is
algebraic of degree~$2$: $H(t)=t^2\sqrt{Q_L(t)}$. The
discriminant $\Delta = 8\kappa S_4$ governs the depth
classification.

\subsection*{6.5.\enspace The four depth classes}

The \emph{shadow depth} $r_{\max}(\cA)$ is the largest degree at
which the shadow obstruction tower carries a nontrivial
invariant, or $\infty$ if the tower never terminates. Three
numerical invariants must be distinguished:
\begin{itemize}[nosep]
\item $p_{\max}(\cA)$: generator OPE pole order,
\item $k_{\max}(\cA)=p_{\max}-1$: collision depth
 (logarithmic absorption),
\item $r_{\max}(\cA)$: shadow depth (independent of
 $p_{\max}$).
\end{itemize}

The standard landscape partitions into four classes
by the cubic shadow $\mathfrak C$ (the degree-$3$ obstruction)
and the quartic contact invariant $\mathfrak Q$ (the
canonical degree-$4$ class of \S6.3). The single-line
dichotomy (Theorem~\ref{thm:single-line-dichotomy}):
$\Delta = 0$ forces tower termination; $\Delta \neq 0$
forces an infinite tower.

\begin{center}
\renewcommand{\arraystretch}{1.3}
\begin{tabular}{lcccl}
\textbf{Class} & $r_{\max}$
 & $\mathfrak C$ & $\mathfrak Q$
 & \textbf{Examples}\\
\hline
Gaussian (G) & $2$ & $0$ & $0$
 & Heisenberg $\cH_k$, free fermion\\
Lie/tree (L) & $3$ & $\neq 0$ & $0$
 & Affine $\widehat{\fg}_k$ (all types)\\
Contact (C) & $4$ & $0$ & $\neq 0$
 & $\beta\gamma$, $bc$ ghosts\\
Mixed (M) & $\infty$ & $\neq 0$ & $\neq 0$
 & Virasoro, $\mathcal W_N$
\end{tabular}
\end{center}

Shadow depth is not Koszulness. All four classes contain
chirally Koszul algebras. Shadow depth classifies
\emph{complexity} of Koszul algebras.

Quintic forcing: the Poisson bracket
$o^{(5)}=\{\mathfrak C,\mathfrak Q\}_H$ forces an infinite
sequence of higher obstructions whenever both $\mathfrak C\neq 0$
and $\mathfrak Q\neq 0$. Finite termination requires
$\mathfrak C=0$ or $\mathfrak Q=0$; when both are nonzero, the
potential is transcendental.

\subsection*{6.6.\enspace The operadic complexity theorem}

The shadow depth equals the $A_\infty$-depth equals the
$L_\infty$-formality level:
\[
r_{\max}(\cA)
\;=\;
A_\infty\text{-depth}(\cA)
\;=\;
L_\infty\text{-formality level}(\gAmod).
\]
The shadow obstruction tower \emph{is} the $L_\infty$-formality
obstruction tower at all degrees
(Theorem~\ref{thm:shadow-formality-identification}). The four
classes G/L/C/M are the four formality types:
$\mathbf{G}$ = free,
$\mathbf{L}$ = Lie,
$\mathbf{C}$ = quadratic non-Lie,
$\mathbf{M}$ = genuinely nonlinear.

\subsection*{6.7.\enspace Genus-two shells}

Genus~$2$ is the first level where all three boundary
strata interact. The MC component decomposes as:
\[
\Theta^{(2)}_\cA
\;=\;
\underbrace{\Theta^{(2)}_{\mathrm{loop}\circ\mathrm{loop}}}_{%
\text{iterated BV}}
\;+\;
\underbrace{\Theta^{(2)}_{\mathrm{sep}\circ\mathrm{loop}}}_{%
\text{clutching $\times$ BV}}
\;+\;
\underbrace{\Theta^{(2)}_{\mathrm{pf}}}_{%
\text{planted-forest}}.
\]
The family-by-family activation:

\begin{center}
\renewcommand{\arraystretch}{1.2}
\begin{tabular}{lcccc}
& $\Theta^{(2)}_{\mathrm{loop}^2}$
& $\Theta^{(2)}_{\mathrm{sep}\circ\mathrm{loop}}$
& $\Theta^{(2)}_{\mathrm{pf}}$
& \textbf{Class}\\
\hline
\textbf{Heisenberg} $\cH_k$
 & $\checkmark$ & --- & --- & G\\
\textbf{Affine} $\widehat{\fg}_k$
 & $\checkmark$ & $\checkmark$ & --- & L\\
\textbf{$\beta\gamma$ system}
 & $\checkmark$ & --- & --- & C\\
\textbf{Virasoro} $\mathrm{Vir}_c$
 & $\checkmark$ & $\checkmark$ & $\checkmark$ & M
\end{tabular}
\end{center}

\noindent
The planted-forest shell is the diagnostic: its presence forces
infinite shadow depth.

\medskip

\noindent\textbf{Shadow growth rate.}\enspace
For class~M algebras, the Taylor coefficients $S_r$ of
$\sqrt{Q_L}$ grow as
$S_r \sim A\cdot\rho(\cA)^r\cdot r^{-5/2}
\cos(r\theta+\varphi)$,
where
$\rho(\cA) = \sqrt{9\alpha^2+2\Delta}\,/\,(2|\kappa|)$
is a continuous invariant. Self-dual Virasoro ($c=13$) has
$\rho\approx 0.467$.

\subsection*{6.8.\enspace The genus loop operator}

The \emph{genus loop operator}
$\Lambda_P\colon\operatorname{Sym}^r(\bar\cA^\vee)_{\mathrm{cyc}}
\to\operatorname{Sym}^{r-2}(\bar\cA^\vee)_{\mathrm{cyc}}$ contracts
two legs of an degree-$r$ cyclic cochain with the propagator:
it raises genus by~$1$ and lowers degree by~$2$, the shadow-level
incarnation of the BV operator $\Delta_{\mathrm{ns}}$.

On the Virasoro primary line, with propagator $P=2/c$ and
quartic contact invariant $Q^{\mathrm{contact}}_{\mathrm{Vir}}
=10/[c(5c{+}22)]$:
\[
\Lambda_P\!\left(\frac{10}{c(5c{+}22)}\cdot x^4\right)
\;=\;
\binom{4}{2}\cdot\frac{2}{c}\cdot
\frac{10}{c(5c{+}22)}\cdot x^2
\;=\;
\frac{120}{c^2(5c{+}22)}\,x^2.
\]
This is the genus-one Hessian correction
$\delta_{H,\mathrm{Vir}}^{(1)}$: a nonlinear genus-$1$
invariant visible only because the shadow tower has depth~$\ge 4$.
The loop ratio
$\rho^{(1)}_{\mathrm{Vir}}
=240/[c^3(5c{+}22)]$
is not determined by $\kappa=c/2$; it detects the quartic
structure through its genus-one projection.

\subsection*{6.9.\enspace Independent sum factorisation}

For a direct sum $L=L_1\oplus L_2$ with vanishing mixed OPE,
all shadow invariants separate: $\kappa$ is additive,
$T^{\mathrm{br}}$ is a direct sum, $\Delta$ is multiplicative,
$R_4^{\mathrm{mod}}$ is additive. The full MC element factorises:
\[
\Theta_{L_1\oplus L_2}
\;=\;
\Theta_{L_1}\hat\otimes 1+1\hat\otimes\Theta_{L_2}.
\]
Every algebra in the standard landscape is either primitive
(Heisenberg, Virasoro, $\beta\gamma$) or built from primitive
pieces by direct sum plus extension. Drinfeld--Sokolov reduction
is the shadow-depth escalator: it sends class~L (finite towers)
to class~M (infinite towers). The nonlinearity of the Dirac
bracket on the Slodowy slice creates the quartic and all higher
shadows.


% ====================================================================
\section*{7.\quad The standard landscape}
% ====================================================================

The machine of Sections~1--6 accepts any chiral algebra. The
seven families below exhaust the standard Lie-theoretic landscape:
each realises a different depth class, curvature sign, and
genus-two shell profile.

\begin{center}
\renewcommand{\arraystretch}{1.2}
\small
\begin{tabular}{llcccccl}
\textbf{Family}
 & \textbf{Generators}
 & $\kappa$
 & $r(z)$
 & $r_{\max}$
 & \textbf{Class}
 & $\kappa+\kappa^!$
 & \textbf{Special}\\
\hline
$\cH_k$
 & $\alpha$ (wt~$1$)
 & $k$
 & $k/z$
 & $2$
 & G
 & $0$
 & Gaussian archetype\\[2pt]
$\widehat{\fg}_k$
 & $J^a$ (wt~$1$)
 & $\frac{(k{+}h^\vee)\dim\fg}{2h^\vee}$
 & $\Omega/z$
 & $3$
 & L
 & $0$
 & FF center at $k{=}{-}h^\vee$\\[4pt]
$\mathrm{Vir}_c$
 & $T$ (wt~$2$)
 & $c/2$
 & $\frac{c/2}{z^3}+\frac{2T}{z}$
 & $\infty$
 & M
 & $13$
 & self-dual at $c{=}13$\\[4pt]
$\cW_3{}_c$
 & $T,W$ (wt~$2,3$)
 & $5c/6$
 & matrix-valued
 & $\infty$
 & M
 & $250/3$
 & cross-channel at $g{\ge}2$\\[4pt]
$\beta\gamma$
 & $\beta,\gamma$ (wt~$\lambda,1{-}\lambda$)
 & $1$
 & $1/z$
 & $4$
 & C
 & $0$
 & contact termination\\[2pt]
$V_\Lambda$
 & $\alpha^i,e^{\lambda\cdot\phi}$
 & $\mathrm{rank}(\Lambda)$
 & Heisenberg
 & $2$
 & G
 & $0$
 & arithmetic from~$q$\\[2pt]
$Y(\fg)$
 & line operators
 & $\mathrm{av}(\Omega/z)$
 & $\Omega/z$
 & $3$
 & L
 & ---
 & YBE $=$ degree-$3$ MC
\end{tabular}
\end{center}

\medskip
\noindent
The landscape census:

\medskip
\noindent\textbf{Heisenberg} (class G).
Single field $\alpha$ of weight~$1$;
OPE $\alpha(z)\alpha(w)\sim k/(z{-}w)^2$. No simple pole:
the bar complex is quadratic, $r_{\max}=2$, and every
higher-genus invariant is determined by $\kappa=k$.
$R$-matrix: $r(z)=k/z$ (scalar, $\Sigma_2$-invariant;
averaging loses nothing). Koszul dual:
$\cH_k^!=(\mathrm{Sym}^{\mathrm{ch}}(V^*),\,m_0=-k\omega)$,
curved. Sewing by Fredholm determinant.

\medskip
\noindent\textbf{Affine Kac--Moody} (class L).
Currents $J^a$ of weight~$1$; OPE with first- and second-order
poles. $R$-matrix: $r(z)=k\Omega/z$, the Casimir tensor of
$\fg$, matrix-valued. Averaging collapses $\Omega$ to
$\kappa(\widehat\fg_k)=(k{+}h^\vee)\dim\fg/(2h^\vee)$. The
matrix structure builds the Yangian $Y(\fg)$; the cubic shadow
is the Lie bracket; the Jacobi identity kills all higher
obstructions.
At critical level $k=-h^\vee$: $\kappa=0$, bar complex flat,
Feigin--Frenkel center
$\mathfrak z(\widehat\fg_{-h^\vee})\cong
\mathrm{Fun}(\mathrm{Op}_{{}^L\fg})$.
Feigin--Frenkel involution $k\leftrightarrow -k-2h^\vee$
is Koszul duality.

\medskip
\noindent\textbf{Virasoro} (class M).
Stress tensor $T$ of weight~$2$; OPE with poles at orders
$4$, $2$, $1$. $R$-matrix:
$r(z)=(c/2)/z^3+2T/z$ (pole orders one less than OPE, by
logarithmic absorption). Self-coupling $T_{(1)}T=2T$ drives
the infinite tower. $Q^{\mathrm{contact}}_{\mathrm{Vir}}
=10/[c(5c{+}22)]$. Quintic
$o^{(5)}=480/[c^2(5c{+}22)]\neq 0$. Koszul dual
$\mathrm{Vir}_{26-c}$; self-dual at $c=13$;
$\kappa+\kappa^!=13$. Critical dimension $c=26$: the dual
curvature $\kappa(\mathrm{Vir}_0)=0$ vanishes.

\medskip
\noindent\textbf{$\mathcal W_3$} (class M).
Two generators $T$ (weight~$2$), $W$ (weight~$3$); the
$WW$~OPE contains the composite quasi-primary
$\Lambda={:}TT{:}-\tfrac{3}{10}\partial^2 T$ with
structure constant $16/(22{+}5c)$. Modular characteristic
$\kappa(\cW_3)=5c/6$. The spectral discriminant is
$\Delta_{\cW_3}(x)=(1-\frac{c}{2}x)(1-\frac{c}{3}x)$.
At genus~$2$, a cross-channel correction from
mixed-propagator graphs appears:
\[
F_2(\cW_3)
\;=\;
\underbrace{\frac{7c}{6912}}_{\kappa\cdot\lambda_2^{\mathrm{FP}}}
\;+\;
\underbrace{\frac{c+204}{16c}}_{\delta F_2},
\]
resolving Open Problem~\ref{op:multi-generator-universality}
negatively: the scalar formula fails for multi-weight algebras
at genus~$\ge 2$.

\medskip
\noindent\textbf{$\beta\gamma$ system} (class C).
Fields $\beta$ (weight~$\lambda$), $\gamma$
(weight~$1{-}\lambda$); OPE
$\beta(z)\gamma(w)\sim 1/(z{-}w)$. Simple pole only, no
double pole. Cubic shadow vanishes (weight mismatch on the
diagonal sector); quartic contact invariant is nonzero but
killed by rank-one abelian rigidity.
$p_{\max}=1$, $k_{\max}=0$, $r_{\max}=4$: the archetypal
witness that shadow depth is independent of OPE pole order.
Koszul dual: $bc$ ghosts.

\medskip
\noindent\textbf{Lattice vertex algebras} (class G).
$V_\Lambda = \cH_k^{\otimes N}\otimes\mathbb C[\Lambda]_\varepsilon$
for an even lattice~$\Lambda$ of rank~$N$.
$\kappa(V_\Lambda)=\mathrm{rank}(\Lambda)$, independent
of~$q$ and~$\varepsilon$. The curvature sees the Heisenberg
factor; the braiding sees $\mathbb C[\Lambda]_\varepsilon$.
Shadow depth~$2$ (Gaussian archetype inherited from
Heisenberg). For $V_\Lambda$ self-dual: the partition function
is an honest modular form. Niemeier discrimination: the bar
complex $(h,\mathrm{rank})$ is a complete invariant for rank-$24$
even self-dual lattice VOAs.

\medskip
\noindent\textbf{Yangians and quantum groups} (class L).
The Yangian $Y(\fg)$ governs the ordered/line-operator sector:
the $R$-matrix $r(z)=k\Omega/z$ is the genus-$0$ binary shadow
of $\Theta^{E_1}_\cA$. The classical Yang--Baxter equation is
Arnold on $\mathrm{FM}_3(\mathbb C)$ evaluated on the Casimir:
three boundary divisors of $\overline C_3(\mathbb C)$, with the
CYBE as $d_{\barB}^2=0$ at degree~$3$, genus~$0$. The RTT
relations $R_{12}R_{13}R_{23}=R_{23}R_{13}R_{12}$ encode the MC
equation at all degrees. The KZ connection
$\nabla_{\mathrm{KZ}}
=d-\hbar\sum_{i<j}\Omega_{ij}\,d\log(z_i{-}z_j)$
is identified with the genus-$0$ shadow connection
$\mathrm{Sh}_{0,n}(\Theta_{\widehat\fg_k})$.

\medskip
\noindent\textbf{Drinfeld--Sokolov reduction as package morphism.}\enspace
$\mathrm{DS}\colon\widehat{\mathfrak{sl}}_3{}_k
\to\cW_3{}_{c(k)}$ extends to a morphism of holographic data.
Shadow depth jumps from~$3$ to~$\infty$: the Dirac bracket on
the arc space $J_\infty(S_f)$ of the Slodowy slice creates the
quartic and all higher shadows. The DS-HPL transfer theorem
makes the chain-level mechanism precise: the $A_\infty$ products
transfer nontrivially, the $r$-matrix transfers with pole-order
shift, but the coproduct is annihilated by a ghost-number
obstruction. All gravitational dynamics resides in the
$r$-matrix.


% ====================================================================
\section*{8.\quad Arithmetic of the shadow obstruction tower}
% ====================================================================

The shadow coefficients $S_r(\cA)$ for the standard families are
rational functions of the central charge. What does number
theory see in these rational numbers? The question has a sharp
answer at genus~$1$, a structural framework at all genera,
and a negative result that delimits the boundary.

\subsection*{8.1.\enspace The shadow Eisenstein theorem}

At genus~$1$, the bar complex over a family of elliptic curves
$\cC_\tau\to\mathfrak H$ produces the amplitude
$\mathrm{Sh}_2^{(1)}(\cA,\tau)=\kappa(\cA)\cdot E_2^*(\tau)$,
where $E_2^*(\tau)=1-24\sum_{n\ge 1}\sigma_1(n)\,q^n$ is the
quasi-modular Eisenstein series (not holomorphic: the
non-holomorphic completion absorbs the anomaly from
$d\log E(z,w)$). The Fourier coefficients $\sigma_1(n)$ are
divisor sums, and their Dirichlet series is a product of
Riemann zeta functions:
\[
D_2(\cA,s)
\;:=\;
\sum_{n\ge 1}\sigma_1(n)\,n^{-s}
\;=\;
\zeta(s)\,\zeta(s{-}1).
\]
The genus-$1$ amplitude Dirichlet series is therefore
$-24\kappa(\cA)\cdot\zeta(s)\cdot\zeta(s{-}1)$: an Eisenstein
$L$-function, with the modular characteristic as the
sole chiral-algebra input.

The shadow $L$-function
$L^{\mathrm{sh}}_\cA(s):=\sum_{r\ge 2}S_r\,r^{-s}$ is a
different object: it extracts the degree-$r$ constant terms of
the shadow tower, not Fourier coefficients of a genus-$1$
amplitude. The distinction is decisive. The shadow
$L$-function is \emph{not} in the Selberg class:
$S_3,\,S_4$ violate the Ramanujan bound for class-M algebras,
and $L^{\mathrm{sh}}_\cA$ has no Euler product (the shadow
coefficients are not multiplicative). The Riemann zeta
function enters through the genus-$1$ Eisenstein channel; the
shadow tower itself carries arithmetic that is not automorphic.

\subsection*{8.2.\enspace The Polyakov correspondence}

Polyakov's programme computes the string partition function
$Z=\sum_g g_s^{2g-2}\int_{\cM_g}[\mathrm{det}'_\zeta\Delta]
^{-d/2}\,d\mu_{\mathrm{WP}}$
by gauge-fixing the worldsheet metric to conformal gauge, paying
the price of a conformal anomaly
$\langle T^a{}_a\rangle=(c/12)\,R$ and a ghost system with
$c_{\mathrm{ghost}}=-26$. Every element of this programme has a
bar-cobar counterpart:
\begin{itemize}[nosep]
\item The conformal anomaly is $\kappa(\cA)$: both measure the
 obstruction to metric-independence at genus~$1$, one by
 functional-integral trace, the other by bar curvature.
\item The critical dimension $c=26$ is the vanishing of the dual
 curvature: $\kappa(\mathrm{Vir}_{26-c})=(26{-}c)/2=0$ when
 $c=26$.
\item The ghost system $bc$ with $c_{\mathrm{ghost}}=-26$ is the
 Koszul dual $\cA^!$ in the BRST pairing: the ghost partition
 function is the complementarity partner.
\item The functional determinant
 $\mathrm{det}'_\zeta\Delta$ on the Laplacian is the Fredholm
 determinant $\det(1-K)$ on the Bergman sewing kernel: the
 Heisenberg one-particle reduction replaces spectral-zeta
 regularisation with a convergent Fredholm series.
\item The genus expansion
 $F_g=\kappa(\cA)\cdot\lambda_g^{\mathrm{FP}}$ (uniform weight)
 is the bar-complex computation of Polyakov's path integral at
 genus~$g$, with the Faber--Pandharipande $\lambda$-class
 replacing Weil--Petersson integration.
\end{itemize}
The four shadow depth classes are the four levels of departure
from the Gaussian path integral: class~G is free ($Z$ is a
determinant), class~L introduces Lie structure (Chern--Simons),
class~C introduces contact terms ($\beta\gamma$ ghosts), and
class~M is full quantum gravity (Polyakov's original problem).

\subsection*{8.3.\enspace The arithmetic packet}

The graph amplitudes at genus~$g\ge 2$ are modular forms for
$\mathrm{Sp}(2g,\mathbb Z)$. The Hecke operators $T_p$ act
on the space of activated amplitudes
$M_\cA=\bigoplus_\chi M_\chi$, decomposed into generalised
eigenspaces. The \emph{arithmetic packet connection}
\[
\nabla^{\mathrm{arith}}_\cA
\;:=\;
d \;-\; \bigoplus_\chi
\frac{\Lambda'_\chi(s)}{\Lambda_\chi(s)}
\bigl(\mathrm{id}_{M_\chi}+N_\chi\bigr)\,ds
\]
is a flat meromorphic connection on $M_\cA\otimes\cO_{\mathbb C}$
over the spectral $s$-line. The structural principle: the
singular divisor depends only on the semisimple part
$M^{\mathrm{ss}}_\cA$ (the arithmetic is semisimple); the
nilpotent terms $N_\chi$ produce unipotent monodromy (the
homotopy defect is unipotent). The Miura splitting for
$\cW_N$: the Heisenberg core is arithmetically inert; all
obstructions beyond Heisenberg lie in a finite defect sector.

For lattice vertex algebras $V_\Lambda$ of rank~$r$, the
\emph{constrained Epstein zeta function}
$\varepsilon^c_s(V_\Lambda)
:=\sideset{}{'}\sum_{\ell\in\Lambda^{\mathrm{bar}}}
\|\ell\|^{-2s}$
restricts the Epstein series to the bar sublattice
$\Lambda^{\mathrm{bar}}\subset\Lambda$. At genus~$2$, the
shadow tautological class computes a central $L$-value via the
B\"ocherer conjecture (Furusawa--Morimoto): the MC element
connects the bar complex to the arithmetic of Siegel modular
forms.


% ====================================================================
\section*{9.\quad The Koszulness programme}
% ====================================================================

\subsection*{9.1.\enspace Scalar saturation}

For a one-parameter algebraic family~$\cA_k$ with rational OPE
coefficients, the universal MC element satisfies: at degree~$2$,
$\Theta_{\cA_k}$ is determined by the single scalar
$\kappa(\cA_k)$. The proof: Whitehead reduction
(vanishing of $H^2(\fg;\mathrm{Sym}^q V)$ for reductive~$\fg$
at generic level) plus algebraic semicontinuity (the exceptional
locus is Zariski-closed, hence empty by the free-field base
point). Scope: every family in the standard landscape at all
non-critical levels, including admissible.

For multi-weight algebras at genus $g\ge 2$, scalar saturation
holds at degree~$2$ but the full MC element receives
cross-channel corrections: the stronger identity
$\Theta_\cA^{\mathrm{min}}=\kappa\cdot\eta\otimes\Lambda$ is
proved on the uniform-weight lane and fails for multi-weight
families (Theorem~\ref{thm:multi-weight-genus-expansion}).

\subsection*{9.2.\enspace The Koszulness meta-theorem}

The Koszul condition admits twelve formulations. Ten are
unconditionally equivalent; one (Lagrangian criterion) is
conditional on perfectness; one ($\mathcal D$-module purity)
is one-directional. For $\cA$ positive-energy and finitely
strongly generated:

\begin{enumerate}[label=\textup{(\roman*)},nosep]
\item $\cA$ is chirally Koszul.
\item PBW spectral sequence collapses at~$E_2$ for every~$g$.
\item Minimal $A_\infty$-model of $H^\bullet(\barB(\cA))$
 is formal: $m_n = 0$ for $n \geq 3$.
\item $\operatorname{Ext}^{p,q}_\cA(\omega_X, \omega_X) = 0$
 for $p \neq q$.
\item Bar-cobar counit
 $\Omega(\barB(\cA)) \to \cA$ is a quasi-isomorphism.
\item Barr--Beck--Lurie comparison is an equivalence.
\item $\int_X \cA$ concentrated in degree~$0$
 for every smooth proper~$X$.
\item $\mathrm{ChirHoch}^*(\cA)$ vanishes outside
 $\{0,1,2\}$.
\item Kac--Shapovalov determinant $\det G_h \neq 0$
 in bar-relevant range.
\item Restriction to every FM boundary stratum acyclic
 outside degree~$0$.
\end{enumerate}

\noindent
Under perfectness/nondegeneracy:

\begin{enumerate}[label=\textup{(\roman*)},nosep,start=11]
\item Lagrangian criterion: $\cM_\cA$ and $\cM_{\cA^!}$ are
 transverse Lagrangians in the $(-1)$-shifted symplectic
 deformation space.
\end{enumerate}

\noindent
And (xii)~$\Rightarrow$~(x), where:

\begin{enumerate}[label=\textup{(\roman*)},nosep,start=12]
\item $\mathcal{D}$-module purity: each $\barB_n(\cA)$ is
 pure of weight~$n$ as a mixed Hodge module. Converse
 open.
\end{enumerate}

\noindent
The shadow-formality identification
(Theorem~\ref{thm:shadow-formality-identification})
identifies the shadow obstruction tower with the
$L_\infty$-formality obstruction tower of the modular
convolution algebra at all degrees. The distinction is precise:
item~(iii) ($A_\infty$-formality of bar cohomology) is a
genus-$0$ condition equivalent to Koszulness; shadow depth
$r_{\max}$ measures the all-genera termination behavior of
$\Theta_\cA$, a transverse invariant invisible to genus~$0$.
All four depth classes G/L/C/M contain chirally Koszul
algebras: the Virasoro algebra has $r_{\max}=\infty$ yet
satisfies item~(iii). Shadow depth classifies complexity
\emph{within} the Koszul world, not Koszulness itself.

\medskip
\noindent\textbf{The proof circuit.}\enspace
The core equivalence chain is
$\textup{(i)} \Leftrightarrow \textup{(ii)}
\Leftrightarrow \textup{(iii)}
\Leftrightarrow \textup{(v)}
\Leftrightarrow \textup{(viii)}$,
proved by PBW concentration $\to$ bar-cobar inversion $\to$
Hochschild polynomial growth $\to$ Koszul-complex acyclicity,
with the $A_\infty$-formality branch via HPL transfer and Keller
classicality. From this core: (iv) by Ext spectral sequence,
(vi) by conservativity + log-FM totalization,
(vii) by the identification $\barB = \int_X$,
(ix) by non-degeneracy forcing PBW injectivity,
(x) by stratum-by-stratum PBW concentration.
The Lagrangian criterion~(xi) uses PTVV derived intersection
theory on the perfectness hypothesis. The $\mathcal{D}$-module
purity implication (xii)$\Rightarrow$(x) uses weight filtration
on the bar complex; the converse would require concentration to
force purity.

\medskip
\noindent\textbf{Additional characterizations.}\enspace
Six further characterizations complement items~(i)--(xii):
shadow-formality at low degree (equivalence at degrees $2,3,4$);
$E_2$-formality of chiral Hochschild cohomology (consequence of
Koszulness); genus-$0$ curve independence (proved independently
by factorization descent); PBW universality (sufficient
condition: every freely strongly generated vertex algebra is
chirally Koszul, covering universal algebras $V_k(\fg)$ at every
level including critical and admissible); tropical Koszulness
(equivalence via the weight spectral sequence on the real
blowup); and bifunctor decomposition (one-directional, converse
open).

For simple quotients $L_k(\fg)$ at admissible levels, the
Koszulness status is rank-dependent: $L_k(\mathfrak{sl}_2)$ is
Koszul at all admissible levels, but at $\mathrm{rk}(\fg)\ge 2$
with denominator $q\ge 3$, Koszulness may fail via a bar
obstruction from the abelian Cartan subalgebra.

\subsection*{9.3.\enspace Chiral Hochschild cohomology (Theorem~H)}

The chiral Hochschild complex
$\mathrm{ChirHoch}^\bullet(\cA)$ is the derived endomorphism
algebra of the identity functor on $\cA$-modules. For
principal $\mathcal W$-algebras at generic level:
\[
\mathrm{ChirHoch}^n(\mathcal W^k(\fg)) = 0
\quad\text{for } n > 2,
\qquad
P(t) = 1 + t^2.
\]
The deformation class in $\mathrm{ChirHoch}^2=\mathbb C$ is the
level deformation $k\mapsto k+\epsilon$. This is a \emph{finite}
invariant (concentrated in three degrees, not a polynomial ring);
contrast with the infinite-dimensional continuous Lie algebra
cohomology $H^*_{\mathrm{GF}}(\mathrm{Diff}(S^1))$, which is
an unrelated invariant. On the Koszul locus,
$\mathrm{ChirHoch}^\bullet(\cA)$ carries a formal $E_2$-algebra
structure.

\subsection*{9.4.\enspace BV/BRST identification}

At genus~$0$, the bar differential is the BRST differential:
$(d_{\barB})_{g=0}=Q_{\mathrm{BRST}}$. Bar-cobar inversion
(Theorem~B) states that BRST cohomology, processed through the
cobar, reconstructs the algebra. At genus~$1$, the
identification acquires curvature: $d^2_{\mathrm{fib}}
=\kappa\cdot\omega_1$ is the one-loop anomaly. At
genus~$g\ge 2$, the chain-level identification is resolved for
classes G, L, and C; for class~M (Virasoro, $\cW_N$), a
coderived reformulation is needed
(Conjecture~\ref{conj:master-bv-brst}).

\subsection*{9.5.\enspace The shadow algebra}

The shadow algebra
$\cA^{\mathrm{sh}}:=H_\bullet(\Defcyc^{\mathrm{mod}}(\cA))$
is a bigraded Lie algebra (graded bracket
$[\kappa,C]\in\cA^{\mathrm{sh}}_{r_1+r_2-2}$; antisymmetric,
Jacobi; the named shadows are projections of MC elements, not
multiplicative generators).
The projections:
\[
\kappa(\cA)=\cA^{\mathrm{sh}}_{2,0},
\qquad
\mathfrak C(\cA)=\cA^{\mathrm{sh}}_{3,0},
\qquad
\mathfrak Q(\cA)=\cA^{\mathrm{sh}}_{4,0},
\qquad
\mathrm{Sh}_r(\cA)=\cA^{\mathrm{sh}}_{r,0}.
\]
The all-degree master equation, read in the shadow algebra:
$\nabla_H(\mathrm{Sh}_r)+o^{(r)}=0$ for all $r\ge 3$. The
shadow algebra is a homotopy invariant of~$\cA$, computed by the
Weil homomorphism of the Chern--Weil dictionary.
